In the first week of the course, due to the absence of two team members, we decided to do some preparation on our own before our first meeting. I chose to delve into the literature to gain a deeper understanding of the fields covered in the "Interaction Technology Innovation" course. Initially, I revisited some research that had made a lasting impression on me. During my undergraduate studies, I participated in a summer school on human-computer interaction organized by my university. Hence, I had previously explored numerous presentations from CHI 2019, including the keynote \textit{Technology Woven In} \cite{poupyrev_2019_chi}. That is why I felt "enchanted" when the original work \textit{Interactive Plant: Botanicus Interacticus} \cite{botanicus} was presented in the first lecture! At the summer school, Professor Yu Chun from Tsinghua University presented a keynote that showcased his research at the time, including HandSee \cite{yu2019handsee}, a design that use a prism to the front camera of a smart phone that enables depth perception of user's hand. I thought it is enchanted because it gives many new interaction possibility to the smartphone that we are so familiar with simply by one prism. There are also works about accessibility, for example, EarTouch \cite{kikuchi2017eartouch} is an interaction design that allow user use their pinna to interact with the screen so that they can hear the phone speaker better in noisy environment. Although they gave me many thoughts, similar ideas are not feasible for us to do. 
In fact, this summer school is the main reason why I chose this Master's program. Reflecting on these enlightening studies brought back many fond memories.

Next, I shifted my focus to gathering broader information by exploring the CHI 2023 - Video Previews playlist on the YouTube ACM SIGCHI channel. I chose to do this because each video in this playlist is about 30 seconds long, making it quite efficient to gain some impressions. I skimmed the related articles if the preview caught my interest. In general, some videos on social connectedness left deep impressions on me. For instance, I was impressed by these two articles:

\begin{enumerate}
    \item Living with Light Touch: An Autoethnography of a Simple Communication Device in Long-Term Use \cite{Livingwithlight}
    \item Message Ritual: A Posthuman Account of Living with Lamp \cite{messageritual}
\end{enumerate}

I found these papers intriguing maybe because I always have a strong interest in social psychology. And the way of using light, a such simple and direct channel, to enhance social connectedness is fascinating. I treated the playlist like a Netflix series, browsing through it and concluding my preparations with some background knowledge.

